%% sec7_1.tex — Section 7.1: Annealing Strategies for Optimization
%% Chapter 7 of "A First Course in Monte Carlo Methods"
%% Sanz-Alonso & Al-Ghattas, University of Chicago

\section{Annealing Strategies for Optimization}
\label{sec:7.1}

In this section, we first use sampling to find the mode of a given target distribution and then
introduce the simulated annealing optimization algorithm to minimize a given objective function.

\subsection{Finding the Mode of a Distribution}
\label{sec:7.1.1}

Let $f$ be the distribution of interest. The mode of $f$ is the set
$\mathcal{M} = \{\xi : f(\xi) \geq f(x) \text{ for all } x \in E\}$
of \textit{global} maxima of $f$.\footnote{We are mainly interested in cases where the mode
comprises a single element $\mathcal{M} = \{\xi\}$, in which case we abuse our terminology and
refer to the element $\xi$ rather than the set $\{\xi\}$ as the mode.}
Naively, one can approximate the mode of $f$ by sampling $f$ and choosing the draw with highest
density:

\begin{enumerate}[label=\arabic*:]
\item Generate $X^{(n)} \sim f$, $\quad 1 \leq n \leq N$.
\item Output $\argmax_{1 \leq n \leq N} f(X^{(n)})$.
\end{enumerate}

This approach is not very efficient. We would like to have a higher density of samples close
to the mode, rather than samples distributed like $f$. This motivates us to modify the
distribution $f$ into a new distribution with the same mode as $f$ but higher density around it.
One way to achieve this goal is to consider
\begin{equation}
  f_\beta(x) \propto f(x)^\beta
  \label{eq:7.1}
\end{equation}
for large values of $\beta > 0$. The new distribution $f_\beta$ is more peaked than $f$ and has
the same mode. We will tacitly assume without further notice that $\int f(x)^\beta\, dx < \infty$,
so that $f(x)^\beta$ can be normalized into a p.d.f.

\begin{example}[Gaussian Target]
\label{ex:7.1}
Consider a Gaussian density $f = \Normal(\mu, \sigma^2)$ with mode $\mu$. In this case
\[
  f_\beta(x) \propto \exp\!\left(-\frac{(x - \mu)^2}{2\sigma^2/\beta}\right),
\]
which shows that $f_\beta(x)$ is a Gaussian density $\Normal(\mu, \sigma^2/\beta)$ with the same
mode $\mu$ as $f$. The larger $\beta$ is chosen, the smaller the variance, and the more
concentrated $f_\beta$ is around its mode.
\end{example}

\begin{example}[Interpretation of $\beta$ as Inverse Temperature]
\label{ex:7.2}
Consider as in Chapter~6 an unnormalized target density of the form
$f(x) \propto \exp(-V(x))$, where $V : \R^d \to \R$ is a given potential function. Then,
$f_\beta(x) \propto \exp(-\beta V(x))$. Following the ideas in Chapter~6, the Langevin equation
\begin{equation}
  dX(t) = -\nabla V(X(t))\, dt + \sqrt{2\beta^{-1}}\, dW(t), \quad 0 \leq t < \infty,
  \label{eq:7.2}
\end{equation}
has $f_\beta$ as invariant distribution under mild assumptions on $V$. The random effect on
the trajectories stemming from the white noise term $dW$ is more pronounced when $\beta^{-1}$
is large. On the other hand, for large $\beta$ the term $\sqrt{2\beta^{-1}}\, dW(t)$ is small
and the Langevin dynamics resembles the deterministic dynamics of the gradient system
$\dot{x} = -\nabla V(x)$. In physical applications, $\beta$ is interpreted as the
\textit{inverse temperature} of the system due to the connection between temperature and
kinetic energy of the system by which particles move faster when the temperature of the system
is higher.
\end{example}

We can use any of the sampling algorithms considered in previous chapters to sample $f_\beta$.
However, sampling $f_\beta$ typically becomes harder for larger $\beta$. Consider for instance
the Random Walk Metropolis Hastings in Algorithm~4.3 with a proposal distribution $g$ symmetric
around $0$. The probability of accepting a move from $X^{(n)}$ to $Z^*$ would be
\begin{equation}
  \min\!\left\{1,\, \frac{f_\beta(Z^*)}{f_\beta(X^{(n)})}\right\}
  = \min\!\left\{1,\, \left(\frac{f(Z^*)}{f(X^{(n)})}\right)^\beta\right\},
  \label{eq:7.3}
\end{equation}
which does not depend on the (usually unknown) normalization constant of $f_\beta$. It is however
difficult to sample from $f_\beta$ for large values of $\beta$. For $\beta \to \infty$ the
probability of accepting a newly proposed $Z^*$ becomes $1$ if $f(Z^*) > f(X^{(n)})$ and $0$
otherwise. For this reason, a typical scenario is that $X^{(n)}$ converges to \textit{local}
extrema of $f$, not necessarily a mode (global extrema).

\begin{example}[Discrete Target and Convergence to Local Maxima]
\label{ex:7.3}
Consider the p.m.f.\ $f(x)$ supported on $\{1, 2, \ldots, 5\}$ defined by
\[
  f(x) =
  \begin{cases}
    0.4 & x = 2, \\
    0.3 & x = 4, \\
    0.1 & x = 1, 3, 5.
  \end{cases}
\]
and shown in Figure~\ref{fig:7.1}. Clearly the mode is $x = 2$. Assume we sample
$f_\beta(x) \propto f(x)^\beta$ with a random walk proposal defined as follows:
\[
  Z^* = X^{(n)} + \xi^*,
\]
where
\begin{align*}
  \Prob(\xi^* = \pm 1) &= 0.5 \quad \text{if } X^{(n)} \in \{2, 3, 4\},\\
  \Prob(\xi^* = -1)    &= 1   \quad \text{if } X^{(n)} = 5,\\
  \Prob(\xi^* = +1)    &= 1   \quad \text{if } X^{(n)} = 1.
\end{align*}
Note that for $\beta \to \infty$ the probability of accepting a move from $4$ to $3$ converges
to $0$ since $f(4) > f(3)$. As the Markov chain can only move from $4$ to $2$ via $3$, it
cannot escape the local mode at $x = 4$ for $\beta \to \infty$.
\end{example}

The above discussion suggests that for large $\beta$ the distribution $f_\beta$ is concentrated
around the mode but is hard to sample from. This motivates considering a time-changing sequence
of target distributions $f_{\beta_n}$ with an increasing sequence $\{\beta_n\}_{n=0}^\infty$ of
inverse temperatures. This is the key idea behind the simulated annealing algorithm we study next.

\subsection{Minimizing an Objective Function}
\label{sec:7.1.2}

We now turn to the problem of finding the global minimum of an objective function
$V : E \to \R$. Equivalently,\footnote{This equivalence makes the implicit assumption that
$\int \exp(-V(x))\, dx < \infty$, so that $\exp(-V(x))$ can be normalized into a p.d.f.}
we seek to find the mode of the distribution
\[
  f(x) \propto \exp\!\bigl(-V(x)\bigr).
\]
As in Subsection~\ref{sec:7.1.1}, we can raise $f$ to different powers to obtain
\[
  f_{\beta_n}(x) \propto \exp\!\bigl(-\beta_n V(x)\bigr),
\]
where $\{\beta_n\}_{n=1}^N$ is a given sequence of inverse temperatures. The simulated annealing
algorithm proceeds as follows:

\begin{figure}[htbp]
\centering
\includegraphics[width=0.75\textwidth]{figures/fig_07_1}
\caption{\textit{An illustration of the target p.m.f.\ in Example~\ref{ex:7.3} along with
$f_\beta(x) \propto f(x)^\beta$ with $\beta = 5, 10$.}}
\label{fig:7.1}
\end{figure}

\begin{algorithm}[H]
\caption{Simulated Annealing}
\label{alg:7.1}
\begin{algorithmic}[1]
\Require Function to minimize $V$, proposal Markov kernel $q(x, z)$, initialization
  $X^{(0)}$, sample size $N$, sequence of inverse temperatures $\{\beta_n\}_{n=1}^N$.
\For{$n = 0, 1, \ldots, N-1$}
  \State Draw $Z^* \sim q(X^{(n)}, \cdot)$.
  \State Update
    \[
      X^{(n+1)} =
      \begin{cases}
        Z^*      & \text{with probability } a := \min\!\left\{1,\,
                   \dfrac{\exp(-\beta_{n+1} V(Z^*))}{\exp(-\beta_{n+1} V(X^{(n)}))}
                   \cdot \dfrac{q(Z^*, X^{(n)})}{q(X^{(n)}, Z^*)}\right\}, \\[6pt]
        X^{(n)} & \text{with probability } 1 - a.
      \end{cases}
    \]
\EndFor
\Ensure Sample $\{X^{(n)}\}_{n=1}^N$ that can be used to minimize $V$ by
  $\argmin_{1 \leq n \leq N} V(X^{(n)})$.
\end{algorithmic}
\end{algorithm}

First and foremost, simulated annealing is a highly versatile optimization algorithm that can
be applied for discrete and continuous optimization without requiring first or second-order
derivatives of the objective. Notice that the accept/reject mechanism agrees with that in the
Metropolis Hastings framework in Algorithm~4.1, taking $f_{\beta_n}$ to be the target
distribution for the $n$-th draw. As in the Metropolis Hastings framework, there is freedom in
the choice of proposal Markov kernel. Since a key feature of simulated annealing is its
potential for derivative-free implementation, random walk proposal kernels that do not require
gradient information on $V$ have been traditionally favored over the Langevin proposal kernels
considered in Chapter~6. If a RWMH with symmetric proposal is used, i.e.\
$Z^* = X^{(n)} + \xi^*$ with the distribution of $\xi^*$ being symmetric around the origin,
then the acceptance probability simplifies to
\[
  a(X^{(n)}, Z^*) = \min\!\left\{1,\, \exp\!\Bigl(-\beta_n\bigl(V(Z^*) - V(X^{(n)})\bigr)\Bigr)\right\},
\]
giving the standard implementation of simulated annealing.

\paragraph{Choice of Inverse Temperatures}
The convergence analysis of simulated annealing is challenging. The algorithm usually converges
to a local minimum, but whether it converges to a global one depends on the choice of inverse
temperatures. Two types of tempering have been widely studied:

\begin{itemize}
\item \textbf{Logarithmic tempering:} For given $\beta_0 > 0$ set
  $\beta_n := \frac{\log(1+n)}{\beta_0}$, $n \geq 1$. The inverse temperatures increase slowly
  enough that global convergence can be established for certain special cases. Simulated
  annealing with logarithmic tempering behaves similarly to an exhaustive search, and so while
  it satisfies nice theoretical guarantees, it is typically slow to converge and has little
  practical relevance.

\item \textbf{Geometric tempering:} $\beta_n = \alpha^n \beta_0$ for some $\beta_0 > 0$ and
  $\alpha > 1$.
\end{itemize}

\begin{center}
\fbox{\parbox{0.92\textwidth}{%
\textbf{Pros and Cons}: Simulated annealing can be an effective discrete or continuous
optimization algorithm when convexity or smoothness assumptions invoked by first and
second-order deterministic optimization algorithms are not satisfied. Importantly, simulated
annealing does not require derivatives and can be applied in high-dimensional settings. The
algorithm converges under mild assumptions when logarithmic tempering is used, but in practice
it is usually implemented with geometric tempering, for which few theoretical guarantees are
available. Additionally, the choice of the temperatures can have a significant impact on the
behavior of the algorithm.
}}
\end{center}
